% =============================================================================
% Chapter 14: Design Choices Explained
% =============================================================================
\chapter{Design Choices Explained}
\label{ch:design-choices}

\begin{objectives}
    \item Understand why TryLM uses RMSNorm over LayerNorm
    \item Learn why RoPE is preferred over learned positions
    \item Grasp why SwiGLU outperforms GELU
    \item Know why pre-normalization aids training
    \item Appreciate weight tying benefits
\end{objectives}

% -----------------------------------------------------------------------------
\section{Overview}
% -----------------------------------------------------------------------------

\TryLM incorporates several modern architectural choices that differ from the original Transformer. Each choice has specific advantages.

\begin{table}[H]
\centering
\begin{tabular}{lll}
\toprule
\textbf{Component} & \textbf{Original} & \textbf{TryLM} \\
\midrule
Normalization & LayerNorm & RMSNorm \\
Position encoding & Learned embeddings & RoPE \\
FFN activation & ReLU/GELU & SwiGLU \\
Norm placement & Post-norm & Pre-norm \\
Output projection & Separate weights & Weight tying \\
Linear bias & Yes & No \\
\bottomrule
\end{tabular}
\caption{Architectural differences from original Transformer}
\end{table}

% -----------------------------------------------------------------------------
\section{RMSNorm vs LayerNorm}
% -----------------------------------------------------------------------------

\subsection{The Difference}

\textbf{LayerNorm}:
\begin{equation}
\text{LN}(x) = \frac{x - \mu}{\sqrt{\sigma^2 + \epsilon}} \cdot \gamma + \beta
\end{equation}

\textbf{RMSNorm}:
\begin{equation}
\text{RMSNorm}(x) = \frac{x}{\sqrt{\frac{1}{d}\sum x_i^2 + \epsilon}} \cdot \gamma
\end{equation}

\subsection{Why RMSNorm?}

\begin{enumerate}
    \item \textbf{Fewer operations}: No mean computation or bias term
    \item \textbf{Empirically equivalent}: No quality loss in practice
    \item \textbf{Used by}: LLaMA, Mistral, and other modern LLMs
\end{enumerate}

\textit{Reference: Zhang \& Sennrich, ``Root Mean Square Layer Normalization'' (2019)}

% -----------------------------------------------------------------------------
\section{RoPE vs Learned Positions}
% -----------------------------------------------------------------------------

\subsection{The Difference}

\textbf{Learned positions}: Add fixed position embeddings to token embeddings.

\textbf{RoPE}: Rotate Q and K vectors based on position.

\subsection{Why RoPE?}

\begin{enumerate}
    \item \textbf{Relative positions}: Attention naturally depends on $(m - n)$, not $m$ and $n$ separately
    \item \textbf{Extrapolation}: Can handle longer sequences than training
    \item \textbf{Efficiency}: Computed with simple sin/cos, no extra parameters
    \item \textbf{Used by}: LLaMA, Mistral, GPT-NeoX
\end{enumerate}

\textit{Reference: Su et al., ``RoFormer: Enhanced Transformer with Rotary Position Embedding'' (2021)}

% -----------------------------------------------------------------------------
\section{SwiGLU vs GELU}
% -----------------------------------------------------------------------------

\subsection{The Difference}

\textbf{GELU}: $\text{GELU}(x) = x \cdot \Phi(x)$ where $\Phi$ is the Gaussian CDF.

\textbf{SwiGLU}: $\text{SwiGLU}(x) = \text{Swish}(xW_g) \odot (xW_u)$

\subsection{Why SwiGLU?}

\begin{enumerate}
    \item \textbf{Gating}: Learned control over information flow
    \item \textbf{Empirical gains}: Consistently outperforms GELU in language modeling
    \item \textbf{Efficiency}: Despite 3 projections vs 2, parameter count is similar (using 2/3 hidden dim)
    \item \textbf{Used by}: LLaMA, PaLM, Mistral
\end{enumerate}

\textit{Reference: Shazeer, ``GLU Variants Improve Transformer'' (2020)}

% -----------------------------------------------------------------------------
\section{Pre-Norm vs Post-Norm}
% -----------------------------------------------------------------------------

\subsection{The Difference}

\textbf{Post-norm}: $x' = \text{Norm}(x + \text{Sublayer}(x))$

\textbf{Pre-norm}: $x' = x + \text{Sublayer}(\text{Norm}(x))$

\subsection{Why Pre-Norm?}

\begin{enumerate}
    \item \textbf{Gradient flow}: Residual path is unobstructed
    \item \textbf{Training stability}: Works with higher learning rates
    \item \textbf{No warmup required}: Post-norm often needs careful warmup
    \item \textbf{Used by}: GPT-2/3, LLaMA, most modern LLMs
\end{enumerate}

\textit{Reference: Xiong et al., ``On Layer Normalization in the Transformer Architecture'' (2020)}

% -----------------------------------------------------------------------------
\section{Weight Tying}
% -----------------------------------------------------------------------------

\subsection{The Technique}

Share weights between input embedding and output projection:
\begin{lstlisting}[style=python]
self.lm_head.weight = self.token_embedding.weight
\end{lstlisting}

\subsection{Why Tie Weights?}

\begin{enumerate}
    \item \textbf{Fewer parameters}: Save $V \times d$ parameters
    \item \textbf{Semantic consistency}: Input and output share representations
    \item \textbf{Regularization}: Constrains the model beneficially
    \item \textbf{Used by}: GPT-2, BERT, most language models
\end{enumerate}

% -----------------------------------------------------------------------------
\section{No Bias Terms}
% -----------------------------------------------------------------------------

\TryLM uses \texttt{bias=False} in linear layers.

\subsection{Why No Bias?}

\begin{enumerate}
    \item \textbf{Fewer parameters}: Small but cumulative savings
    \item \textbf{Simplicity}: Easier to analyze
    \item \textbf{RMSNorm compatibility}: Bias less necessary with normalization
    \item \textbf{Empirical}: No quality loss observed
\end{enumerate}

% -----------------------------------------------------------------------------
\section{Summary}
% -----------------------------------------------------------------------------

\begin{summary}
    \item \textbf{RMSNorm}: Simpler than LayerNorm, equally effective
    \item \textbf{RoPE}: Encodes relative positions, enables extrapolation
    \item \textbf{SwiGLU}: Gated activation outperforms GELU
    \item \textbf{Pre-norm}: Better gradient flow, more stable training
    \item \textbf{Weight tying}: Shares input/output embeddings, saves parameters
    \item These choices follow modern best practices from LLaMA/Mistral
\end{summary}
